\chapter{Literature and Research Rationale}
\label{ch:literature}
\section{Physics-Informed Flow Approximation}
Raissi et al.\ formulate PINNs by penalizing differential-equation residuals
alongside observed or prescribed information \cite{raissi2019}. The broader
physics-informed learning framework includes several ways to combine
physical constraints and learned approximations \cite{karniadakis2021}.
For this project, the relevant distinction is between a network trained on
PDE and boundary information and a regression model trained directly on
interior CFD labels. The current flow model uses the former design, and its
FEM interior samples are reserved for numerical comparison.

Optimization behavior limits this approach. Gradient imbalance among loss
terms and unequal learning rates of solution components have been studied
in PINNs \cite{wang2021,wang2022}. These results motivate scaling the
momentum and continuity residuals, retaining fixed monitoring samples, and
reporting boundary and flux errors separately. They do not establish that
any particular optimizer or architecture will be accurate in the proposed
curved-channel problem.

\section{Independent Flow and Particle References}
Mixed finite elements provide a classical route to incompressible flow.
The Taylor--Hood construction uses a higher-order velocity space than
pressure space \cite{taylor1973}; the present independent reference uses
P2 velocity and P1 pressure with scikit-fem \cite{gustafsson2020}.
Independence here means that this solver is assembled without PINN
predictions. Its own nonlinear, mesh, and continuity checks remain necessary.
It is a project-specific CFD reference, not a claim that commercial CFD or
an experimental benchmark was used.

The full small-sphere equation of motion includes terms beyond instantaneous
Stokes drag \cite{maxey1983}. The present model deliberately uses a
heavy-particle simplification. This is appropriate only within a qualified
regime and cannot represent every particle--fluid interaction. In particular,
geometric wall contact does not supply the hydrodynamic wall corrections
that may become important at small clearances.

\section{Verification, Validation, and the Specific Gap}
Numerical verification addresses consistency with the stated equations;
physical validation concerns agreement with suitable external physical
evidence \cite{roache1998}. Analytical Poiseuille flow and a manufactured
three-dimensional solution are strong implementation checks. Agreement with
a refined FEM solution tests numerical consistency for the same model.
Neither substitutes for experiment when making device-level predictions.

\begin{table}[!htbp]
\centering\small
\caption{Roles of the main methodological foundations.}
\label{tab:literature}
\begin{tabularx}{\textwidth}{@{}p{.23\textwidth}YY@{}}
\toprule Foundation & Role in this proposal & Limit of inference\\\midrule
PINNs \cite{raissi2019,karniadakis2021} & Differentiable flow approximation constrained by physics & Training loss does not certify trajectory accuracy.\\
Training analysis \cite{wang2021,wang2022} & Motivation for loss scaling and separate diagnostics & Remedies must be evaluated on this channel family.\\
Taylor--Hood FEM \cite{taylor1973} & Independent velocity/pressure reference & Discretization and geometry errors remain.\\
Small-sphere dynamics \cite{maxey1983} & Basis for the drag-only heavy-particle model & Other force terms and confinement are omitted.\\
Numerical verification \cite{roache1998} & Reference hierarchy and error reporting & Numerical agreement is not experimental validation.\\
\bottomrule
\end{tabularx}
\end{table}

The proposed gap is local to the current workflow: its global field metrics
and particle-outcome metrics do not yet support the same level of confidence.
The study will investigate that discrepancy and determine a defensible
operating envelope. This is a focused narrative rationale based on the cited
foundational work and repository evidence; it is not presented as an
exhaustive or systematic survey of all particle-transport PINNs.
